% -----------------------------------------------------------------------------
% Goetsu sioji (月絶書) PUA glyphs — goetsusioji.ttf
% Load after preamble/fonts-and-ruby.tex and preamble/jyutcitzi-mode.tex.
%
% PUA code points share Unicode space with Jyutcitzi.  Do not call \jczOff here:
% that switches xeCJK to BabelStone Han, which lacks PUA glyphs, and with
% fallback disabled they become tofu.  Instead, wrap content in \goetsufont only.
%
%   \begin{goetsu}…\end{goetsu}
%   \goetsuinline{…}
% -----------------------------------------------------------------------------

\makeatletter

\newif\ifgoetsu@defined
\goetsu@definedfalse

\newcommand\goetsu@define{%
  \newfontfamily{\goetsufont}{goetsusioji}[Path=\goetsu@path, Extension=.ttf]%
  \newCJKfontfamily{\goetsuCJK}{goetsusioji}[Path=\goetsu@path, Extension=.ttf]%
  \goetsu@definedtrue
}

% Real goetsusioji.ttf is ~13 MB (internal name: Source Han Serif CN).
% A ~24 MB file named goetsusioji.ttf but containing Jyutcitzi is the wrong copy.

\IfFileExists{../fonts/goetsusioji.ttf}{%
  \def\goetsu@path{../fonts/}%
  \goetsu@define
}{%
  \IfFileExists{fonts/goetsusioji.ttf}{%
    \def\goetsu@path{fonts/}%
    \goetsu@define
  }{%
    \IfFileExists{../../fonts/goetsusioji.ttf}{%
      \def\goetsu@path{../../fonts/}%
      \goetsu@define
    }{%
      \PackageWarningNoLine{goetsu}{%
        goetsusioji.ttf not found (tried ../fonts/, fonts/, ../../fonts/);%
        goetsu environment will pass text through unchanged%
      }%
    }%
  }%
}

\newcommand{\goetsu@typeset}[1]{%
  \begingroup
  \goetsufont{\ignorespaces#1\unskip}%
  \endgroup
}

\newcommand{\goetsuinline}[1]{%
  \ifgoetsu@defined
    \goetsu@typeset{#1}%
  \else
    #1%
  \fi
}

\ExplSyntaxOn
\NewDocumentEnvironment { goetsu } { +b }
  {
    \ifgoetsu@defined
      \goetsu@typeset {#1}
    \else
      #1
    \fi
  }
  { }
\ExplSyntaxOff

\makeatother
